\hypertarget{appendix-5d4252a42886}{}
\section{Auto decltype and reference collapsing}\label{appendix-5d4252a42886-auto-decltype-and-reference-collapsing}

\textbf{Curriculum depth:} Introduced · \textbf{First seen:} C008 · \textbf{Latest development:} C008

\subsection{Definition and mechanism}\label{appendix-5d4252a42886-definition-and-mechanism}

Ordinary by-value auto drops top-level const/reference information; auto\& preserves referent qualification. decltype(name) reports the declared type, while decltype((expression)) follows category. Collapsing yields an lvalue reference if either participating reference is \&, otherwise \&\&.

\subsection{Prerequisites}\label{appendix-5d4252a42886-prerequisites}

\begin{itemize}
\tightlist
\item
  \hyperref[appendix-1121b30eaeb0]{Value categories and materialization} --- Expression category determines viable binding and whether temporary identity is needed.
\item
  \hyperref[appendix-d6d562ca9fbb]{References and aliases} --- The initial alias and its referent determine later binding and deduction behavior.
\end{itemize}

\subsection{Used by}\label{appendix-5d4252a42886-used-by}

\begin{itemize}
\tightlist
\item
  \hyperref[appendix-0104b704fb11]{Move casts and forwarding references}
\end{itemize}

\subsection{Recall and next depth}\label{appendix-5d4252a42886-recall-and-next-depth}

Explain why decltype(rr) is T\&\& but decltype((rr)) is T\&. Pending: constrained deduction, structured bindings and return-type deduction. Written coverage does not establish learner mastery.

\subsection{Source and navigation}\label{appendix-5d4252a42886-source-and-navigation}

\hyperref[01-foundations-i-ch003-an-expression-has-a-type-and-a-category]{Teaching explanation and examples} · \hyperref[appendix-fd090d982a15-m008---parts-113-126]{Exact coverage} · \href{../02-Foundations-II/main.pdf\#nameddest=appendix-4ad47bdd269a}{Dependency Map (02-Foundations-II, appendix-4ad47bdd269a)}
